\documentclass[12pt]{article}
\usepackage{amsmath,amssymb,amsfonts,amsthm}
\usepackage[margin=1in]{geometry}

\begin{document}

\begin{center}
{\Large\bfseries Chapter 9, Problem 15}\\[6pt]
Byron \& Fuller, \textit{Mathematics of Classical and Quantum Physics}
\end{center}

\bigskip

\textbf{Problem.}

(a) Using Problem~9.14(a), show $N(x,y;\lambda) = k(x,y)D(\lambda) + \lambda\int k(x,t)N(t,y;\lambda)\,dt$ and similarly from 9.14(b).

(b) Define $R(x,y;\lambda) = N(x,y;\lambda)/D(\lambda)$. Show that $(I - \lambda K)(I + \lambda R) = I$ and $(I + \lambda R)(I - \lambda K) = I$, identifying $R$ as the resolvent of $K$.

\bigskip
\textbf{Solution.}

\subsection*{(a) The two fundamental equations}

From Problem~9.14(a), expanding $\Delta_n$ along the first row and summing the series, we established:
\[
N(x,y;\lambda) = k(x,y)\,D(\lambda) + \lambda\int_a^b k(x,t)\,N(t,y;\lambda)\,dt. \tag{I}
\]

From Problem~9.14(b), expanding along the first column:
\[
N(x,y;\lambda) = k(x,y)\,D(\lambda) + \lambda\int_a^b N(x,t;\lambda)\,k(t,y)\,dt. \tag{II}
\]

\subsection*{(b) Resolvent identification}

Define $R(x,y;\lambda) = N(x,y;\lambda)/D(\lambda)$, assuming $D(\lambda) \neq 0$.

Dividing equation~(I) by $D(\lambda)$:
\[
R(x,y;\lambda) = k(x,y) + \lambda\int k(x,t)\,R(t,y;\lambda)\,dt.
\]

In operator notation, $R = K + \lambda KR$, i.e., $R - \lambda KR = K$, so:
\[
(I - \lambda K)R = K \quad \Longrightarrow \quad R = (I - \lambda K)^{-1}K.
\]

Now check: $(I - \lambda K)(I + \lambda R) = I - \lambda K + \lambda R - \lambda^2 KR$. Using $R = K + \lambda KR$:
\begin{align}
\lambda R &= \lambda K + \lambda^2 KR, \\
\lambda R - \lambda^2 KR &= \lambda K.
\end{align}
Therefore:
\[
(I - \lambda K)(I + \lambda R) = I - \lambda K + \lambda K = I. \quad \checkmark
\]

Similarly, dividing equation~(II) by $D(\lambda)$:
\[
R(x,y;\lambda) = k(x,y) + \lambda\int R(x,t;\lambda)\,k(t,y)\,dt,
\]
i.e., $R = K + \lambda RK$, so $(I + \lambda R)(I - \lambda K) = I + \lambda R - \lambda K - \lambda^2 RK$. Using $\lambda R = \lambda K + \lambda^2 RK$:
\[
(I + \lambda R)(I - \lambda K) = I + \lambda K + \lambda^2 RK - \lambda K - \lambda^2 RK = I. \quad \checkmark
\]

Therefore $I + \lambda R(\lambda)$ is the two-sided inverse of $I - \lambda K$, and $R(\lambda)$ is the \textbf{resolvent} of $K$.

The integral equation $f = g + \lambda Kf$ can thus be solved as:
\[
f = (I - \lambda K)^{-1}g = g + \lambda R(\lambda)g = g + \lambda\int_a^b R(x,y;\lambda)\,g(y)\,dy,
\]
valid for all $\lambda$ such that $D(\lambda) \neq 0$ (i.e., $\lambda$ is not a singular value of $K$). \hfill $\blacksquare$

\end{document}
